\input{estilo}
% Propuesta: arriba, productor, cola y consumidor simétricos, con un RPC de
% cada lado; el receive es una llamada del consumidor a la cola (la flecha va
% al revés de como la dibujó la pizarra, cosa que la clase aclara). Abajo, las
% formas de recibir como cuatro paneles iguales: blocking, polling, push y
% notify. La pizarra las listaba; los diagramas salen de lo que dice la clase.
\begin{document}
\begin{figura}\begin{tikzpicture}
  % --- el esquema
  \node[caja, minimum width=24mm, minimum height=13mm] (p) at (0,0) {productor};
  \node[cola, minimum height=30mm] (q) at (5.0,0) {};
  \node[caja, minimum width=24mm, minimum height=13mm] (c) at (10.0,0) {consumidor};
  \draw[msg] (p.east) -- node[cod, above] {push(\var{m})} (q.west);
  \draw[msg] ($(c.west)+(0,0.2)$) -- node[cod, above] {receive()} ($(q.east)+(0,0.2)$);
  \draw[respuesta] ($(q.east)+(0,-0.2)$) -- node[etiq, below] {\var{m}} ($(c.west)+(0,-0.2)$);
  \foreach \a/\b in {1.2/3.5, 6.5/8.8}
    \draw[llave] (\a,0.75) -- (\b,0.75) node[midway, above=6pt, etiq, text=suave] {RPC};
  \draw[llave] (6.5,-0.75) -- (3.5,-0.75) node[midway, below=6pt, etiq, text=suave] {cola de mensajes};

  % --- las formas de recibir: cuatro paneles iguales
  \def\yt{-2.3} \def\yb{-6.0}
  \newcommand\panel[2]{% centro, título
    \node[etiq] at (#1-0.8,\yt) {cola}; \draw[vida] (#1-0.8,\yt-0.25) -- (#1-0.8,\yb);
    \node[etiq] at (#1+0.8,\yt) {consumidor}; \draw[vida] (#1+0.8,\yt-0.25) -- (#1+0.8,\yb);
    \node at (#1,\yb-0.45) {#2};}
  \newcommand\pide[2]{\draw[msg] (#1+0.8,#2) -- (#1-0.8,#2-0.2);}
  \newcommand\da[3]{\draw[respuesta] (#1-0.8,#2) -- node[etiq, above, sloped] {#3} (#1+0.8,#2-0.2);}
  \newcommand\llega[2]{\node[evento=acento] at (#1-0.8,#2) {};
    \node[etiq, text=acento, anchor=east] at (#1-0.9,#2) {llega \var{m}};}

  % blocking: el consumidor se queda esperando
  \panel{0.05}{blocking}
  \pide{0.05}{-2.9}
  \fill[suave!40] (0.77,-3.1) rectangle (0.93,-4.55);
  \node[etiq, text=suave, anchor=west] at (0.97,-3.8) {espera};
  \llega{0.05}{-4.35}
  \da{0.05}{-4.45}{\var{m}}

  % polling: pregunta, no hay nada, vuelve a preguntar
  \panel{3.35}{polling}
  \foreach \y/\r in {-2.8/vacío, -3.85/vacío, -4.9/\var{m}} {
    \pide{3.35}{\y}
    \da{3.35}{\y-0.3}{\r}
  }
  \llega{3.35}{-4.75}

  % push: la cola se lo manda
  \panel{6.65}{push}
  \llega{6.65}{-3.3}
  \da{6.65}{-3.4}{\var{m}}

  % notify: la cola avisa y el consumidor lo va a buscar
  \panel{9.95}{notify}
  \llega{9.95}{-2.7}
  \draw[msg] (9.15,-3.0) -- node[etiq, above, sloped] {hay uno} (10.75,-3.2);
  \pide{9.95}{-3.75}
  \da{9.95}{-4.05}{\var{m}}
\end{tikzpicture}\end{figura}
\end{document}
